\section{Beadwork is a ported thread-walk, not a strand link}

The beadwork claim-site is not a strand crossing but a ported rigid unit with a
thread walk through holes.  This subgraph separates the port carrier from the
thread walk: bead ports are $\charge$-typed coherence cells, and a threading
programme is a directed walk on the port graph.

\begin{theorem}[The bead port is a charge cell]\label{thm:bead-portgraph-carrier}\lean{Obsmetrologia.AperiodicOrder.Tiling.over_occupation_inconsistency}\leanok\uses{thm:crossing-charge, thm:status-correspondence, thm:condensate-is-top}
A bead with its two principal ports
$(p_h^-,p_h^+)$ carries a four-corner datum --- the pair (entered? , exited?) --- landing in
\texttt{Charge} through \texttt{Articulating.FourCorner.toCorner} exactly as the weave crossing does
(\cref{thm:crossing-charge}). A doubly-traversed bead (thread passes twice, as in
right-angle weave) is the held $\top_k$ of a bead; an un-entered hole is $\bot_k$
(\texttt{condensate}; \texttt{over\_occupation\_inconsistency}). The beadwork stitch alphabet is a
charge carrier, not the binary \texttt{Bool} of \texttt{warpAlt}. The ported-cell carrier now
exists (\texttt{Obsmetrologia.AperiodicOrder.Tiling.PortGraph}); the residue is the multi-hole transition system and
the four adjacencies.
\end{theorem}

\begin{theorem}[The four bead adjacencies have a finite carrier]\label{thm:bead-four-adjacency-carrier}\lean{Obsmetrologia.AperiodicOrder.Crystal.BeadAdjacency, Obsmetrologia.AperiodicOrder.Crystal.BeadAdjacency.tuple_eq, Obsmetrologia.AperiodicOrder.Crystal.Beads.threadAdjacency, Obsmetrologia.AperiodicOrder.Crystal.Beads.contactAdjacency, Obsmetrologia.AperiodicOrder.Crystal.Beads.cellAdjacency, Obsmetrologia.AperiodicOrder.Crystal.Beads.visibleAdjacency, Obsmetrologia.AperiodicOrder.Crystal.Beads.adjacencyTuple_eq}\leanok\uses{thm:bead-portgraph-carrier}
The bead layer now carries four named finite graphs on the same lenticule
carrier: thread, contact, cell, and visible adjacency.  This proves the
multi-adjacency carrier itself.  It does not yet prove that the four relations
are pairwise distinct or that the visible graph is a sofic image of the thread
walk.
\end{theorem}

\begin{theorem}[Thread adjacency is not contact adjacency]\label{thm:bead-thread-walk-distinct}\notready\uses{thm:bead-four-adjacency-carrier, thm:charge-coloured-sofic}
\emph{(Open frontier --- stated, not formalized.)} The four adjacency graphs of
\cref{thm:bead-four-adjacency-carrier} should be distinct in non-strandlike
beadwork.  The thread
relation is the directed walk $\Gamma=(b_{i_1},\dots,b_{i_N})$ with repeats allowed; the
visible relation is the $1$-block charge factor (\cref{thm:charge-coloured-sofic}) read off
the finished surface. The claim: the visible-surface sofic shift is the forgetful image of
the thread walk, and the two coincide only for strand-like stitches (loom beadwork). A
ported port graph $P_B=(V_P,E_{\text{hole}},E_{\text{thread}},E_{\text{contact}})$ is the
genuine carrier; the \texttt{Path} development sees only the thread walk.
\end{theorem}

\begin{theorem}[Graph threading is a settling closed walk]\label{thm:bead-threading-closed-walk}\notready\uses{thm:bead-thread-walk-distinct, thm:interlacement-is-settle}
\emph{(Open frontier --- stated, not formalized.)} The Demaine--Kirkpatrick--Lin
graph-threading problem --- route one thread through the holes so that pulling it tight
realizes a target bead graph --- is a settling on the port carrier: the tightened
realization is the past-closure of the threading walk, and a legal threading exists iff the
closed-walk cover settles to the target (the bead reading of
\cref{thm:interlacement-is-settle}). Minimal thread length is the least-fixed-point cover;
an illegal port turn is a point where the settle fails to absorb. Requires the bead-port
carrier of \cref{thm:bead-portgraph-carrier} mapped to a \texttt{SettleFrame}.
\end{theorem}
